\documentclass[12pt, a4paper]{article}
\usepackage[utf8]{inputenc}
\usepackage{geometry}
\usepackage{hyperref}
\usepackage{enumitem}

\geometry{top=1in, bottom=1in, left=1in, right=1in}

\title{\textbf{Electric Motorike Digital Twin}}
\author{
  \Large Makau Abigael  (ENM221-0052/2022) \\
  \Large Koros Patrobas  (ENM221-0074/2022)
}
\date{}

\begin{document}

\maketitle

\section {Introduction and Background}
\large 
The rapid electrification of the transport sector , particularly adoption of electric motorbikes in emerging markets has expose a critical engineering : high cost, time latency and safety risks associated with traditional physical prototyping. Iterative testing of electric powertrain components and battery management systems traditionally requires fully assembled mortocycles. This approach is vulnerable to expensive failures such as motor burnout, battery thermal runaway,power-controller burnout and mechanical stress fractures when subjected to untested real world conditions.
\paragraph{}This project proposesthe development of a Hardware-in-the-loop (HIL) digitl twin architecture for electric motorbikes that simulates the dynamic behavior of the powertrain and battery system under various load conditions and real-world scenarios. 
\paragraph{}The Mechatronic testing board will allow for physical simulation of motor dynamics, battery drain and spatial loading characteristics. this enables engineers to stress test physical components and validate control algorithms against virtual terrain and rider payloads safely and cost-effectively. The digital twin dashboard will provide real-time telemetry, predictive maintenance alerts and performance analytics, enabling data-driven design decisions and optimization of electric motorbike systems. 
\section{Main Objective of the Study}
To design and develop a fully functional electric Motorbike Digital Twin
capable of real-time performance monitoring and predictive maintenance.

\section{Specific Objectives}
\begin{enumerate}
    \item To formulate a mathematical model of an electric motorbike's powertrain, caculating spatial kinetics, aerodynamic drag and variable loading parameters.
    \item to design and assemble the physical testing boarsd with integrated sensors, actuators and microcontrollers to replicate motor and battery behavior.
    \item to impleent a real-time, closed-loop control system that establishes a two-way data link, ensuring the controller dynamically adjusts physical motor torque based on virtual terrain calculated by the simulation model.
    \item To implement a digital twin dashboard that visualizes simulated powertrain metrics and evaluates system response under different load profiles.
  \end{enumerate}

\section{Expected Outcomes}
\begin{enumerate}
    \item  A functional hardware benchtop model capable of mimicking electric motorbike dynamics and load profiles.
    \item A digital twin dashboard providing real-time visual feedback on motor speed, torque characteristics, and battery depletion rates.
    \item A safe, repeatable platform for evaluating electric powertrain efficiency and stress limits without risking full vehicle hardware.
\end{enumerate}

\section{Justification}
Xxx \\
Xxx \\
Xxx.

\section{General Research Methodology}
The project will be executed in the following phases:
\begin{enumerate}[label=\alph*)]
    \item \textbf{Hardware Testbench Design:} Constructing the physical testing board framework, incorporating variable load elements, a test motor, and simulated power sources.
    \item \textbf{Electronic Circuit and Sensor Integration:} Embedding microcontrollers , current sensors, voltage dividers, and speed encoders onto the board.
    \item \textbf{Software and Dashboard Development:} Creating the digital twin visualization interface to map incoming telemetry to virtual gauges and graphs.
    \item \textbf{System Testing and Evaluation:} Running different loading and motor profiles on the test board to verify real-time mirroring and accuracy against the physical simulation.
\end{enumerate}

\section{Partnerships or Collaborations}
(If any)

\section{Conclusion}
The Automatic xxx.

\vspace{1em}
\noindent\textbf{Contacts:}
\begin{itemize}
    \item Makau Abigael  
    \\\textbf{Email:} \textit{musengyabigael@gmail.com}
   \\ \textbf{Tel:} 0797122059
    \item Koros Patrobas  
    \\\textbf{Email:} \textit{}
   \\ \textbf{Tel:} 0745216899
\end{itemize}

\end{document}
